% ============================================================
%  Section 2 – Background
% ============================================================

\begin{definition}[GGM Tree]
\label{def:ggm}
Let $\PRG : \{0,1\}^{\lam} \to \{0,1\}^{2\lam}$ be a length-doubling pseudorandom
generator, decomposed as $\PRG(x) = (\PRGl(x),\, \PRGr(x))$, where $\PRGl$ and
$\PRGr$ each output $\lam$ bits.  Given a root seed
$\seed_r \in \{0,1\}^{\lam}$ and a depth $d$, the \emph{GGM tree} is a complete
binary tree of depth $d$ in which:
\begin{itemize}
  \item the root holds $\seed_r$;
  \item each internal node $v$ storing seed $s$ generates its left child as
        $\PRGl(s)$ and its right child as $\PRGr(s)$;
  \item the $2^{d}$ leaves are the \PRG\ outputs at depth $d$.
\end{itemize}
\end{definition}

The crucial property exploited by ZK signature schemes is the
\emph{punctured-PRF property}: given all leaf seeds \emph{except} leaf~$i$, it is
computationally infeasible to recover the seed at leaf~$i$, provided $\PRG$ is
a secure pseudorandom generator.  This makes the \GGM\ tree suitable for hiding
one or more seeds while revealing the rest.

\paragraph{Incomplete \GGM\ trees.}
When the number of required leaves $t$ is not a power of $2$, a scheme may use an
\emph{incomplete} tree: leaves appear at multiple levels, and internal nodes near
the bottom may have fewer than two children or act as leaves themselves.  This
reduces memory and computation, but the additional index-tracking logic introduces
new fault surfaces, as detailed in Section~\ref{sec:theoretical}.

% ============================================================
\subsection{LESS-v2}
\label{sec:less}
% ============================================================

\paragraph{Context.}
LESS (Linear Equivalence Signature Scheme)~\cite{LESS-spec} uses a Fiat--Shamir
transformation over a 3-pass Sigma protocol whose security rests on the Canonical
Form Linear Equivalence Problem (CF-LEP).  During signing, the signer must
generate $t$ independent ephemeral monomial matrices
$\mathbf{Q}_{e_0}, \ldots, \mathbf{Q}_{e_{t-1}} \in \Mon_n(q)$.

\paragraph{Tree type.}
LESS-v2 uses an \emph{incomplete} \GGM\ tree with $T = \lceil \log_2 t \rceil$
levels.  Leaves are distributed across the bottom two levels; the index offsets
are tracked via the arrays $\mathsf{npl}$, $\mathsf{ll}$, $\mathsf{off}$, and
$\mathsf{lsi}$.

\paragraph{Flag tree.}
After Fiat--Shamir challenge generation, LESS-v2 constructs an auxiliary
\emph{flag tree} of the same shape as the seed tree.  It is a labelling
$\calF : V(\calT) \to \{0,1\}$ where $\calF(v) = 0$ marks a publishable
node and $\calF(v) = 1$ marks a hidden node.  The flag tree is built in three
sequential phases:
\begin{enumerate}
  \item \textbf{Initialisation.}  Every node is set to $0$ except the root,
        which is set to $1$.
  \item \textbf{Leaf labelling} ($\LabelLeaves$).  For each leaf $i$, set
        \begin{equation*}
          \calF(\mathrm{leaf}_i) =
          \begin{cases} 1 & \text{if } b[i] \neq 0, \\ 0 & \text{otherwise,} \end{cases}
        \end{equation*}
        where $\mathbf{b} = (b[0], \ldots, b[t-1])$ is the Fiat--Shamir challenge
        vector of weight $w$ over the alphabet $\{0, \ldots, s-1\}$.
  \item \textbf{Upward propagation} ($\ComputeSeeds$).  For each internal node $v$
        with children $\ell$ and $r$:
        \begin{equation*}
          \calF(v) = \calF(\ell) \vee \calF(r).
        \end{equation*}
        This ensures the minimum set of nodes is published to reconstruct all
        revealed leaf seeds.
\end{enumerate}

\paragraph{Secret-extraction opportunity.}
For every round $i$ with $b[i] \neq 0$, the signer publishes
$\CosetRep\!\bigl(\pi_{b[i]} \circ \tilde{\pi}^{*}_i\bigr)$
while keeping $\eseed_i$ hidden.  If an attacker obtains both values
simultaneously, Theorem~2 of~\cite{FaultToForge2025} shows that the secret
permutation $\pi_{b[i]}$ is recoverable.  The flag tree is the
\emph{sole enforcement mechanism} that prevents this dual revelation.

\begin{example}
Figure~\ref{fig:less-tree} illustrates an incomplete \GGM\ tree for $t = 6$ with
challenge $\mathbf{b} = (0,1,0,1,0,1)$, so
$\calH = \{1,3,5\}$ and $\calR = \{0,2,4\}$.
After all three phases, nodes $s_3$, $s_7$, and $s_9$ satisfy
$\calF(v)=0$ with $\calF(\mathsf{parent}(v))=1$ and thus enter $\tpath$,
giving the verifier access to $\eseed_0$, $\eseed_2$, and $\eseed_4$.
\end{example}

\begin{figure}[tb]
\centering
\begin{tikzpicture}[x=0.95cm, y=1.15cm]
  % ---- nodes (final state after all 3 phases) ----
  % Level 0
  \node[hidnode]  (s0)  at (3.5, 0)   {$s_0$};
  % Level 1
  \node[hidnode]  (s1)  at (1.5,-1)   {$s_1$};
  \node[hidnode]  (s2)  at (5.5,-1)   {$s_2$};
  % Level 2  (s3,s4 are leaves; s5,s6 are internal)
  \node[pathnode] (s3)  at (0.5,-2)   {$s_3$};
  \node[hidnode]  (s4)  at (2.5,-2)   {$s_4$};
  \node[hidnode]  (s5)  at (4.5,-2)   {$s_5$};
  \node[hidnode]  (s6)  at (6.75,-2)   {$s_6$};
  % Level 3  (leaves)
  \node[pathnode] (s7)  at (3.75,-3)  {$s_7$};
  \node[hidnode]  (s8)  at (5.25,-3)  {$s_8$};
  \node[pathnode] (s9)  at (6.00,-3)  {$s_9$};
  \node[hidnode]  (s10) at (7.25,-3)  {$s_{10}$};
  % ---- ephemeral-seed labels below leaves ----
  \node[leaflbl] at (0.5,-2.55)  {$\eseed_0$};
  \node[leaflbl] at (2.5,-2.55)  {$\eseed_1$};
  \node[leaflbl] at (3.75,-3.55) {$\eseed_2$};
  \node[leaflbl] at (5.25,-3.55) {$\eseed_3$};
  \node[leaflbl] at (5.75,-3.55) {$\eseed_4$};
  \node[leaflbl] at (7.25,-3.55) {$\eseed_5$};
  % ---- edges ----
  \draw[treeedge] (s0)--(s1); \draw[treeedge] (s0)--(s2);
  \draw[treeedge] (s1)--(s3); \draw[treeedge] (s1)--(s4);
  \draw[treeedge] (s2)--(s5); \draw[treeedge] (s2)--(s6);
  \draw[treeedge] (s5)--(s7); \draw[treeedge] (s5)--(s8);
  \draw[treeedge] (s6)--(s9); \draw[treeedge] (s6)--(s10);
  % ---- "incomplete" annotation ----
  \node[font=\tiny,gray] at (3.5,-2.55) {leaf level~2};
  % ---- legend ----
  \matrix[draw=gray!40, rounded corners=2pt, inner sep=4pt,
          above left=0.2cm and -0.8cm of s1,
          matrix of nodes, nodes={font=\tiny, inner sep=2pt},
          column sep=4pt, row sep=1pt] {
    \node[hidnode,  minimum size=8pt]{}; & \node{$\calF=1$ (hidden)};  \\
    \node[pathnode, minimum size=8pt]{}; & \node{$\calF=0$, in $\tpath$}; \\
  };
\end{tikzpicture}
\caption{LESS-v2 incomplete \GGM\ flag tree, $t=6$, after all three construction
         phases.  Challenge $\mathbf{b}=(0,1,0,1,0,1)$ hides rounds
         $\calH=\{1,3,5\}$.  Path nodes (blue, thick border) give the verifier
         $\eseed_0$, $\eseed_2$, $\eseed_4$.  Injecting a fault on any red node
         exposes the corresponding hidden seed alongside its response, yielding a
         dual revelation.}
\label{fig:less-tree}
\end{figure}


Table~\ref{tab:less-params} lists the parameter sets of LESS.

\begin{table}[tb]
\centering
\caption{LESS-v2 parameter sets~\cite{LESS-spec}.}
\label{tab:less-params}
\begin{tabular}{@{}lcccccc@{}}
\toprule
Parameter set   & $n$ & $k$ & $q$ & $t$ & $w$ & $s$ \\
\midrule
LESS-252-45     & 252 & 126 & 127 &  45 & 34 & 8 \\
LESS-252-68     & 252 & 126 & 127 &  68 & 42 & 4 \\
LESS-252-192    & 252 & 126 & 127 & 192 & 36 & 2 \\
LESS-400-102    & 400 & 200 & 127 & 102 & 61 & 4 \\
LESS-400-220    & 400 & 200 & 127 & 220 & 68 & 2 \\
LESS-548-137    & 548 & 274 & 127 & 137 & 79 & 4 \\
LESS-548-345    & 548 & 274 & 127 & 345 & 75 & 2 \\
\bottomrule
\end{tabular}
\end{table}

% ============================================================
\subsection{MEDS}
\label{sec:meds}
% ============================================================

\paragraph{Context.}
MEDS (Matrix Equivalence Digital Signature)~\cite{ChouNPRRST23} is a code-based
signature scheme whose security rests on the Matrix Code Equivalence Problem
(MCEP).  Like LESS, it applies the Fiat--Shamir transform to a 3-pass Sigma
protocol between a prover holding a secret equivalence
$(A, B) \in \GL_k(q) \times \GL_n(q)$ and a verifier.

\paragraph{Tree type.}
MEDS uses a \emph{complete} binary \GGM\ seed tree to generate $t$ ephemeral
seeds.  Each leaf seed $\eseed_i$ expands into an ephemeral matrix pair.  The
binary challenge vector $\mathbf{h} = (h_0, \ldots, h_{t-1}) \in \{0,1\}^t$
determines which seeds are revealed ($h_i = 0$) and which are hidden ($h_i = 1$).

\paragraph{Reference tree.}
MEDS uses a \emph{Reference Tree} (also called $\SeedTree$) with a
$\SeedTreePaths$ algorithm.  This traverses the tree bottom-up and returns the
minimal set of internal nodes sufficient to reconstruct all revealed leaf seeds,
while no hidden seed $\eseed_i$ (with $h_i = 1$) is derivable.  The three
abstract phases map as follows:
\begin{enumerate}
  \item \textbf{Expansion.}
        $\BuildSeeds(\seed_r, \salt)$ expands from the root to all $t$ leaves via
        $\PRG(\text{node} \,\|\, \salt \,\|\, \text{index}) \to (\text{left},\,\text{right})$.
  \item \textbf{Selection.}
        $\SeedTreePaths(\calT, \mathbf{h})$ marks and collects the publishing nodes.
  \item \textbf{Reconstruction (verifier).}
        $\ReconstructSeeds(\tpath, \mathbf{h}, \salt)$ re-expands only from the
        published nodes.
\end{enumerate}

\paragraph{Secret-extraction opportunity.}
For round $i$ with $h_i = 1$, the signer reveals its ZK response.  If an attacker
obtains both the response and $\eseed_i$, the secret equivalence pair $(A, B)$ is
recoverable by linear algebra.  The Reference Tree is the sole enforcement mechanism
preventing this dual revelation.

\begin{example}
Figure~\ref{fig:meds-tree} shows a complete \GGM\ tree for $t = 4$ with
challenge $\mathbf{h} = (0,1,0,1)$, so $\calH = \{1,3\}$ and $\calR = \{0,2\}$.
Nodes $s_3$ and $s_5$ enter $\tpath$.
\end{example}

\begin{figure}[tb]
\centering
\begin{tikzpicture}[x=1.1cm, y=1.15cm]
  % Level 0
  \node[hidnode]  (s0) at (2,  0)  {$s_0$};
  % Level 1
  \node[hidnode]  (s1) at (0.5,-1) {$s_1$};
  \node[hidnode]  (s2) at (3.5,-1) {$s_2$};
  % Level 2 (leaves)
  \node[pathnode] (s3) at (-0.5,-2){$s_3$};
  \node[hidnode]  (s4) at ( 1.5,-2){$s_4$};
  \node[pathnode] (s5) at ( 2.5,-2){$s_5$};
  \node[hidnode]  (s6) at ( 4.5,-2){$s_6$};
  % Ephemeral seed labels
  \node[leaflbl] at (-0.5,-2.55){$\eseed_0$};
  \node[leaflbl] at ( 1.5,-2.55){$\eseed_1$};
  \node[leaflbl] at ( 2.5,-2.55){$\eseed_2$};
  \node[leaflbl] at ( 4.5,-2.55){$\eseed_3$};
  % Edges
  \draw[treeedge] (s0)--(s1); \draw[treeedge] (s0)--(s2);
  \draw[treeedge] (s1)--(s3); \draw[treeedge] (s1)--(s4);
  \draw[treeedge] (s2)--(s5); \draw[treeedge] (s2)--(s6);
  % Legend
  \matrix[draw=gray!40, rounded corners=2pt, inner sep=4pt,
          above right=0.6cm and -1cm of s1,
          matrix of nodes, nodes={font=\tiny, inner sep=2pt},
          column sep=4pt, row sep=1pt] {
    \node[hidnode,  minimum size=8pt]{}; & \node{$\calF=1$};  \\
    \node[pathnode, minimum size=8pt]{}; & \node{in $\tpath$}; \\
  };
\end{tikzpicture}
\caption{MEDS complete \GGM\ flag tree, $t=4$.
         Challenge $\mathbf{h}=(0,1,0,1)$ hides $\calH=\{1,3\}$.
         Path nodes $s_3$ and $s_5$ (blue) give the verifier $\eseed_0$ and
         $\eseed_2$.  A fault on $s_4$ or $s_6$ exposes the corresponding hidden
         seed alongside its ZK response.}
\label{fig:meds-tree}
\end{figure}

\begin{table}[tb]
\centering
\caption{MEDS parameter sets~\cite{ChouNPRRST23}.}
\label{tab:meds-params}
\begin{tabular}{@{}lccccc@{}}
\toprule
Parameter set & $n$ & $k$ & $q$ & $t$ & $s$ \\
\midrule
MEDS-9923b    & 14 &  7 & 509 & 192 & 2 \\
MEDS-13220b   & 22 & 11 & 127 & 192 & 2 \\
MEDS-41711b   & 30 & 15 & 509 & 392 & 2 \\
\bottomrule
\end{tabular}
\end{table}

% ============================================================
\subsection{CROSS}
\label{sec:cross}
% ============================================================

\paragraph{Context.}
CROSS (Codes and Restricted Objects Signature Scheme)~\cite{CROSSspec} is based on
the Restricted Syndrome Decoding Problem (R-SDP)~\cite{BaldiBPSWW24} and applies the
Fiat--Shamir transform to a 5-pass ($q_2$-identification) protocol.
Like MEDS, CROSS uses a \emph{binary} fixed-weight challenge vector
$\mathbf{c} = (c_0, \ldots, c_{t-1}) \in \{0,1\}^t$ of weight $w$,
so the hidden and revealed index sets are
$\calH = \{i : c_i = 1\}$ and $\calR = \{i : c_i = 0\}$,
identical in structure to MEDS.

\paragraph{Tree type.}
CROSS uses a complete \GGM\ seed tree.  Round~1 of CROSS uses the same
complete binary tree as LESS-v1.  The Round~2 specification~\cite{CROSSspec}
adopts a combined seed-and-Merkle-tree structure, but the seed-publishing
component remains structurally equivalent to the $\SeedTreePaths$ mechanism
of MEDS (Section~\ref{sec:meds}).

\paragraph{Publishing mechanism.}
Flag propagation follows exactly the same three-phase bottom-up rule as
MEDS: initialise all flags to~$0$ except the root, label leaves from the
binary challenge, propagate upward by disjunction.  The published path gives
the verifier access to all $\eseed_i$ with $c_i = 0$, while every $\eseed_i$
with $c_i = 1$ remains hidden.  The structural identity with MEDS means
that every fault location identified for MEDS (Section~\ref{sec:meds}) applies
to CROSS without modification.

\paragraph{Secret-extraction opportunity.}
For round $i$ with $c_i = 1$, the signer reveals a ZK response encoding the
relationship between $\eseed_i$ and the long-term secret R-SDP solution.
If an attacker obtains both the response and $\eseed_i$, the secret is
recoverable by the linear-algebra argument of ZKFault~\cite{ZKFault},
which demonstrated full key recovery from a single fault for all CROSS
parameter sets.

\begin{example}
Figure~\ref{fig:cross-tree} shows a complete \GGM\ tree for $t = 4$ with
binary challenge $\mathbf{c} = (0,1,0,1)$, so $\calH = \{1,3\}$ and
$\calR = \{0,2\}$.  Nodes $s_3$ and $s_5$ enter $\tpath$.
The tree topology and the flag-propagation logic are identical to the MEDS
example of Figure~\ref{fig:meds-tree}; only the underlying hard problem
and response format differ.  A fault on $s_4$ or $s_6$ exposes the
corresponding hidden seed alongside its R-SDP response, enabling recovery
of the long-term secret.
\end{example}

\begin{figure}[tb]
\centering
\begin{tikzpicture}[x=1.1cm, y=1.15cm]
  % Level 0
  \node[hidnode]  (s0) at (2,  0)  {$s_0$};
  % Level 1
  \node[hidnode]  (s1) at (0.5,-1) {$s_1$};
  \node[hidnode]  (s2) at (3.5,-1) {$s_2$};
  % Level 2 (leaves)
  \node[pathnode] (s3) at (-0.5,-2){$s_3$};
  \node[hidnode]  (s4) at ( 1.5,-2){$s_4$};
  \node[pathnode] (s5) at ( 2.5,-2){$s_5$};
  \node[hidnode]  (s6) at ( 4.5,-2){$s_6$};
  % Ephemeral-seed labels + challenge values
  \node[leaflbl] at (-0.5,-2.6){$\eseed_0,\; c_0{=}0$};
  \node[leaflbl] at ( 1.5,-2.6){$\eseed_1,\; c_1{=}1$};
  \node[leaflbl] at ( 2.5,-2.6){$\eseed_2,\; c_2{=}0$};
  \node[leaflbl] at ( 4.5,-2.6){$\eseed_3,\; c_3{=}1$};
  % Edges
  \draw[treeedge] (s0)--(s1); \draw[treeedge] (s0)--(s2);
  \draw[treeedge] (s1)--(s3); \draw[treeedge] (s1)--(s4);
  \draw[treeedge] (s2)--(s5); \draw[treeedge] (s2)--(s6);
  % Legend
  \matrix[draw=gray!40, rounded corners=2pt, inner sep=4pt,
          above right=0.6cm and -1cm of s1,
          matrix of nodes, nodes={font=\tiny, inner sep=2pt},
          column sep=4pt, row sep=1pt] {
    \node[hidnode,  minimum size=8pt]{}; & \node{$\calF=1$ (hidden)};  \\
    \node[pathnode, minimum size=8pt]{}; & \node{$\calF=0$, in $\tpath$}; \\
  };
\end{tikzpicture}
\caption{CROSS complete \GGM\ flag tree, $t=4$, binary challenge
         $\mathbf{c}=(0,1,0,1)$.  The tree structure is identical to the MEDS
         example (Figure~\ref{fig:meds-tree}): path nodes $s_3$ and $s_5$ (blue)
         give the verifier $\eseed_0$ and $\eseed_2$.  A fault on $s_4$ or
         $s_6$ exposes the corresponding hidden seed alongside its R-SDP
         response, enabling full key recovery as shown in ZKFault~\cite{ZKFault}.}
\label{fig:cross-tree}
\end{figure}

% ============================================================
\subsection{MPCitH / VOLEitH Family (FAEST, Mirath, SDitH, RYDE)}
\label{sec:mpcith}
% ============================================================

\paragraph{Context.}
MPC-in-the-Head (MPCitH) schemes use the \GGM tree in a more
elaborate but structurally analogous way.  They generate $N$ virtual
\emph{party shares} for each of $\tau$ parallel repetitions.  For repetition $e$,
the Fiat--Shamir challenge selects one hidden party index
$i^*_e \in \{0, \ldots, N-1\}$.  The \GGM\ tree lets the signer reveal all $N-1$
non-hidden shares using $O(\log N)$ published intermediate seeds.

\paragraph{Tree types.}
\begin{itemize}
  \item \textbf{Classic MPCitH (Mirath, RYDE).}
        A forest of $\tau$ independent complete binary trees, each with $N$
        leaves.  All roots are derived from a single master seed via $\XOF$.
  \item \textbf{VOLEitH (FAEST, SDitH v2).}
        A single large \GGM\ tree with $\tau \times N$ leaves, as described
        in~\cite{OneTree}.  Child-node generation uses $\PRG = \text{AES-CTR}$
        (FAEST) or SHAKE (SDitH).
\end{itemize}

The all-but-one vector commitment (AVC / BAVC) mechanism is the MPCitH counterpart
of the flag tree.  For each repetition $e$, the signer publishes the
\emph{co-path}: the sequence of sibling nodes along the path from the hidden leaf
$i^*_e$ to the root.  The verifier can re-expand all $N-1$ non-hidden leaf seeds
from the co-path alone.  This is structurally the \GGM-path mechanism with
exactly one hidden leaf per repetition, and the same fault surfaces (index
arithmetic, loop increments, conditional branch) apply.

In LESS, MEDS, and CROSS, the secret is a permutation or matrix equivalence
embedded in the hidden seed's ephemeral commitment.  In MPCitH, the hidden seed
generates the prover's witness share.  In both cases the invariant is identical:
the signer must never simultaneously reveal both the hidden leaf seed \emph{and}
the corresponding ZK response for the same round.
